% Options for packages loaded elsewhere
\PassOptionsToPackage{unicode}{hyperref}
\PassOptionsToPackage{hyphens}{url}
%
\documentclass[
]{article}
\usepackage{amsmath,amssymb}
\usepackage{lmodern}
\usepackage{iftex}
\ifPDFTeX
  \usepackage[T1]{fontenc}
  \usepackage[utf8]{inputenc}
  \usepackage{textcomp} % provide euro and other symbols
\else % if luatex or xetex
  \usepackage{unicode-math}
  \defaultfontfeatures{Scale=MatchLowercase}
  \defaultfontfeatures[\rmfamily]{Ligatures=TeX,Scale=1}
\fi
% Use upquote if available, for straight quotes in verbatim environments
\IfFileExists{upquote.sty}{\usepackage{upquote}}{}
\IfFileExists{microtype.sty}{% use microtype if available
  \usepackage[]{microtype}
  \UseMicrotypeSet[protrusion]{basicmath} % disable protrusion for tt fonts
}{}
\makeatletter
\@ifundefined{KOMAClassName}{% if non-KOMA class
  \IfFileExists{parskip.sty}{%
    \usepackage{parskip}
  }{% else
    \setlength{\parindent}{0pt}
    \setlength{\parskip}{6pt plus 2pt minus 1pt}}
}{% if KOMA class
  \KOMAoptions{parskip=half}}
\makeatother
\usepackage{xcolor}
\IfFileExists{xurl.sty}{\usepackage{xurl}}{} % add URL line breaks if available
\IfFileExists{bookmark.sty}{\usepackage{bookmark}}{\usepackage{hyperref}}
\hypersetup{
  pdftitle={Exercise 2},
  pdfauthor={David Jorgensen},
  hidelinks,
  pdfcreator={LaTeX via pandoc}}
\urlstyle{same} % disable monospaced font for URLs
\usepackage[margin=1in]{geometry}
\usepackage{color}
\usepackage{fancyvrb}
\newcommand{\VerbBar}{|}
\newcommand{\VERB}{\Verb[commandchars=\\\{\}]}
\DefineVerbatimEnvironment{Highlighting}{Verbatim}{commandchars=\\\{\}}
% Add ',fontsize=\small' for more characters per line
\usepackage{framed}
\definecolor{shadecolor}{RGB}{248,248,248}
\newenvironment{Shaded}{\begin{snugshade}}{\end{snugshade}}
\newcommand{\AlertTok}[1]{\textcolor[rgb]{0.94,0.16,0.16}{#1}}
\newcommand{\AnnotationTok}[1]{\textcolor[rgb]{0.56,0.35,0.01}{\textbf{\textit{#1}}}}
\newcommand{\AttributeTok}[1]{\textcolor[rgb]{0.77,0.63,0.00}{#1}}
\newcommand{\BaseNTok}[1]{\textcolor[rgb]{0.00,0.00,0.81}{#1}}
\newcommand{\BuiltInTok}[1]{#1}
\newcommand{\CharTok}[1]{\textcolor[rgb]{0.31,0.60,0.02}{#1}}
\newcommand{\CommentTok}[1]{\textcolor[rgb]{0.56,0.35,0.01}{\textit{#1}}}
\newcommand{\CommentVarTok}[1]{\textcolor[rgb]{0.56,0.35,0.01}{\textbf{\textit{#1}}}}
\newcommand{\ConstantTok}[1]{\textcolor[rgb]{0.00,0.00,0.00}{#1}}
\newcommand{\ControlFlowTok}[1]{\textcolor[rgb]{0.13,0.29,0.53}{\textbf{#1}}}
\newcommand{\DataTypeTok}[1]{\textcolor[rgb]{0.13,0.29,0.53}{#1}}
\newcommand{\DecValTok}[1]{\textcolor[rgb]{0.00,0.00,0.81}{#1}}
\newcommand{\DocumentationTok}[1]{\textcolor[rgb]{0.56,0.35,0.01}{\textbf{\textit{#1}}}}
\newcommand{\ErrorTok}[1]{\textcolor[rgb]{0.64,0.00,0.00}{\textbf{#1}}}
\newcommand{\ExtensionTok}[1]{#1}
\newcommand{\FloatTok}[1]{\textcolor[rgb]{0.00,0.00,0.81}{#1}}
\newcommand{\FunctionTok}[1]{\textcolor[rgb]{0.00,0.00,0.00}{#1}}
\newcommand{\ImportTok}[1]{#1}
\newcommand{\InformationTok}[1]{\textcolor[rgb]{0.56,0.35,0.01}{\textbf{\textit{#1}}}}
\newcommand{\KeywordTok}[1]{\textcolor[rgb]{0.13,0.29,0.53}{\textbf{#1}}}
\newcommand{\NormalTok}[1]{#1}
\newcommand{\OperatorTok}[1]{\textcolor[rgb]{0.81,0.36,0.00}{\textbf{#1}}}
\newcommand{\OtherTok}[1]{\textcolor[rgb]{0.56,0.35,0.01}{#1}}
\newcommand{\PreprocessorTok}[1]{\textcolor[rgb]{0.56,0.35,0.01}{\textit{#1}}}
\newcommand{\RegionMarkerTok}[1]{#1}
\newcommand{\SpecialCharTok}[1]{\textcolor[rgb]{0.00,0.00,0.00}{#1}}
\newcommand{\SpecialStringTok}[1]{\textcolor[rgb]{0.31,0.60,0.02}{#1}}
\newcommand{\StringTok}[1]{\textcolor[rgb]{0.31,0.60,0.02}{#1}}
\newcommand{\VariableTok}[1]{\textcolor[rgb]{0.00,0.00,0.00}{#1}}
\newcommand{\VerbatimStringTok}[1]{\textcolor[rgb]{0.31,0.60,0.02}{#1}}
\newcommand{\WarningTok}[1]{\textcolor[rgb]{0.56,0.35,0.01}{\textbf{\textit{#1}}}}
\usepackage{graphicx}
\makeatletter
\def\maxwidth{\ifdim\Gin@nat@width>\linewidth\linewidth\else\Gin@nat@width\fi}
\def\maxheight{\ifdim\Gin@nat@height>\textheight\textheight\else\Gin@nat@height\fi}
\makeatother
% Scale images if necessary, so that they will not overflow the page
% margins by default, and it is still possible to overwrite the defaults
% using explicit options in \includegraphics[width, height, ...]{}
\setkeys{Gin}{width=\maxwidth,height=\maxheight,keepaspectratio}
% Set default figure placement to htbp
\makeatletter
\def\fps@figure{htbp}
\makeatother
\setlength{\emergencystretch}{3em} % prevent overfull lines
\providecommand{\tightlist}{%
  \setlength{\itemsep}{0pt}\setlength{\parskip}{0pt}}
\setcounter{secnumdepth}{-\maxdimen} % remove section numbering
\ifLuaTeX
  \usepackage{selnolig}  % disable illegal ligatures
\fi

\title{Exercise 2}
\author{David Jorgensen}
\date{2025-12-14}

\begin{document}
\maketitle

\hypertarget{data-prep}{%
\section{Data Prep}\label{data-prep}}

\begin{Shaded}
\begin{Highlighting}[]
\CommentTok{\# load CSV into data frame \textasciigrave{}ledger\textasciigrave{}}
\NormalTok{ledger }\OtherTok{\textless{}{-}} \FunctionTok{read.csv}\NormalTok{(}\StringTok{"sales\_data\_sample.csv"}\NormalTok{)}
\CommentTok{\# set ledger$OrderDate to Date class}
\NormalTok{ledger}\SpecialCharTok{$}\NormalTok{OrderDate }\OtherTok{\textless{}{-}} \FunctionTok{as.Date}\NormalTok{(ledger}\SpecialCharTok{$}\NormalTok{OrderDate, }\AttributeTok{format=}\StringTok{"\%m/\%d/\%Y"}\NormalTok{)}
\CommentTok{\# create new column Month extracting month from OrderDate}
\NormalTok{ledger}\SpecialCharTok{$}\NormalTok{Month }\OtherTok{\textless{}{-}} \FunctionTok{as.numeric}\NormalTok{(}\FunctionTok{format}\NormalTok{(ledger}\SpecialCharTok{$}\NormalTok{OrderDate, }\StringTok{"\%m"}\NormalTok{))}
\CommentTok{\# check the structure of the data}
\FunctionTok{str}\NormalTok{(ledger)}
\end{Highlighting}
\end{Shaded}

\begin{verbatim}
## 'data.frame':    495 obs. of  8 variables:
##  $ OrderID   : int  1 2 3 4 5 6 7 8 9 10 ...
##  $ OrderDate : Date, format: "2024-01-12" "2025-08-27" ...
##  $ Region    : chr  "South" "South" "North" "West" ...
##  $ Product   : chr  "Product_D" "Product_A" "Product_A" "Product_A" ...
##  $ Quantity  : int  7 15 11 8 7 19 11 11 4 8 ...
##  $ UnitPrice : num  34.37 8.59 2 3.09 24.05 ...
##  $ TotalSales: num  60.1 30.1 34 23.4 707.2 ...
##  $ Month     : num  1 8 5 4 9 10 8 4 4 5 ...
\end{verbatim}

\begin{Shaded}
\begin{Highlighting}[]
\CommentTok{\# check the class of the data}
\FunctionTok{class}\NormalTok{(ledger)}
\end{Highlighting}
\end{Shaded}

\begin{verbatim}
## [1] "data.frame"
\end{verbatim}

\begin{Shaded}
\begin{Highlighting}[]
\CommentTok{\# assign ledger to sales\_data in global environment}
\NormalTok{sales\_data }\OtherTok{\textless{}\textless{}{-}}\NormalTok{ ledger}
\end{Highlighting}
\end{Shaded}

\hypertarget{charts}{%
\section{Charts}\label{charts}}

\hypertarget{bar}{%
\subsection{Bar}\label{bar}}

\begin{Shaded}
\begin{Highlighting}[]
\CommentTok{\# Bar graph showing product and the sum of sales for each product}
\CommentTok{\# aggregate sales by product}
\NormalTok{sales\_by\_product }\OtherTok{\textless{}{-}} \FunctionTok{aggregate}\NormalTok{(sales\_data}\SpecialCharTok{$}\NormalTok{TotalSales, }\AttributeTok{by=}\FunctionTok{list}\NormalTok{(}\AttributeTok{product=}\NormalTok{sales\_data}\SpecialCharTok{$}\NormalTok{Product), }\AttributeTok{FUN=}\NormalTok{sum)}
\CommentTok{\# set column names}
\FunctionTok{colnames}\NormalTok{(sales\_by\_product) }\OtherTok{\textless{}{-}} \FunctionTok{c}\NormalTok{(}\StringTok{"product"}\NormalTok{, }\StringTok{"sales"}\NormalTok{)}
\CommentTok{\# create barplot}
\FunctionTok{barplot}\NormalTok{(sales\_by\_product}\SpecialCharTok{$}\NormalTok{sales, }\AttributeTok{names.arg=}\NormalTok{sales\_by\_product}\SpecialCharTok{$}\NormalTok{product,}
\AttributeTok{main=}\StringTok{"Sales by Product"}\NormalTok{, }\AttributeTok{xlab=}\StringTok{"Product"}\NormalTok{, }\AttributeTok{ylab=}\StringTok{"Total Sales"}\NormalTok{)}
\end{Highlighting}
\end{Shaded}

\includegraphics{djorgensen-520-exercise2_files/figure-latex/bar-1.pdf}

\hypertarget{pie}{%
\subsection{Pie}\label{pie}}

\begin{Shaded}
\begin{Highlighting}[]
\CommentTok{\# Pie chart showing sum of sales by region}
\NormalTok{sales\_by\_region }\OtherTok{\textless{}{-}} \FunctionTok{aggregate}\NormalTok{(sales\_data}\SpecialCharTok{$}\NormalTok{TotalSales, }\AttributeTok{by=}\FunctionTok{list}\NormalTok{(}\AttributeTok{region=}\NormalTok{sales\_data}\SpecialCharTok{$}\NormalTok{Region), }\AttributeTok{FUN=}\NormalTok{sum)}
\CommentTok{\# set column names}
\FunctionTok{colnames}\NormalTok{(sales\_by\_region) }\OtherTok{\textless{}{-}} \FunctionTok{c}\NormalTok{(}\StringTok{"region"}\NormalTok{, }\StringTok{"sales"}\NormalTok{)}
\CommentTok{\# create pie chart}
\FunctionTok{pie}\NormalTok{(sales\_by\_region}\SpecialCharTok{$}\NormalTok{sales, }\AttributeTok{labels=}\NormalTok{sales\_by\_region}\SpecialCharTok{$}\NormalTok{region,}
\AttributeTok{main=}\StringTok{"Sales by Region"}\NormalTok{)}
\end{Highlighting}
\end{Shaded}

\includegraphics{djorgensen-520-exercise2_files/figure-latex/pie-1.pdf}

\hypertarget{box-plot}{%
\subsection{Box Plot}\label{box-plot}}

\begin{Shaded}
\begin{Highlighting}[]
\CommentTok{\# Box plots for average sales by region and average quantities by region}
\CommentTok{\# boxplot of TotalSales by Region}
\FunctionTok{boxplot}\NormalTok{(TotalSales }\SpecialCharTok{\textasciitilde{}}\NormalTok{ Region, }\AttributeTok{data=}\NormalTok{sales\_data,}
\AttributeTok{main=}\StringTok{"Total Sales by Region"}\NormalTok{, }\AttributeTok{xlab=}\StringTok{"Region"}\NormalTok{, }\AttributeTok{ylab=}\StringTok{"Total Sales"}\NormalTok{)}
\end{Highlighting}
\end{Shaded}

\includegraphics{djorgensen-520-exercise2_files/figure-latex/boxplot-1.pdf}

\begin{Shaded}
\begin{Highlighting}[]
\CommentTok{\# boxplot of Quantity by Region}
\FunctionTok{boxplot}\NormalTok{(Quantity }\SpecialCharTok{\textasciitilde{}}\NormalTok{ Region, }\AttributeTok{data=}\NormalTok{sales\_data,}
\AttributeTok{main=}\StringTok{"Quantity by Region"}\NormalTok{, }\AttributeTok{xlab=}\StringTok{"Region"}\NormalTok{, }\AttributeTok{ylab=}\StringTok{"Quantity"}\NormalTok{)}
\end{Highlighting}
\end{Shaded}

\includegraphics{djorgensen-520-exercise2_files/figure-latex/boxplot-2.pdf}

\hypertarget{histogram}{%
\subsection{Histogram}\label{histogram}}

Because this data is categorical, a histogram is not appropriate.
Instead, we can create a bar plot to show the count of orders by
product.

\begin{Shaded}
\begin{Highlighting}[]
\CommentTok{\# Barplot of the count of orders by product.}
\CommentTok{\# create table of counts by product}
\NormalTok{orders\_by\_product }\OtherTok{\textless{}{-}} \FunctionTok{table}\NormalTok{(sales\_data}\SpecialCharTok{$}\NormalTok{Product)}
\CommentTok{\# create barplot}
\FunctionTok{barplot}\NormalTok{(orders\_by\_product,}
\AttributeTok{main=}\StringTok{"Count of Orders by Product"}\NormalTok{, }\AttributeTok{xlab=}\StringTok{"Product"}\NormalTok{, }\AttributeTok{ylab=}\StringTok{"Number of Orders"}\NormalTok{)}
\end{Highlighting}
\end{Shaded}

\includegraphics{djorgensen-520-exercise2_files/figure-latex/histogram-1.pdf}

\hypertarget{radial-plot}{%
\subsection{Radial Plot}\label{radial-plot}}

\begin{Shaded}
\begin{Highlighting}[]
\CommentTok{\# A radial plot showing total sales (for all products) by month.}
\FunctionTok{library}\NormalTok{(ggplot2)}

\CommentTok{\# Aggregate total sales by month}
\NormalTok{sales\_by\_month }\OtherTok{\textless{}{-}} \FunctionTok{aggregate}\NormalTok{(sales\_data}\SpecialCharTok{$}\NormalTok{TotalSales, }\AttributeTok{by=}\FunctionTok{list}\NormalTok{(}\AttributeTok{month=}\NormalTok{sales\_data}\SpecialCharTok{$}\NormalTok{Month), }\AttributeTok{FUN=}\NormalTok{sum)}
\CommentTok{\# set column names}
\FunctionTok{colnames}\NormalTok{(sales\_by\_month) }\OtherTok{\textless{}{-}} \FunctionTok{c}\NormalTok{(}\StringTok{"month"}\NormalTok{, }\StringTok{"sales"}\NormalTok{)}
\CommentTok{\# Create quarter column}
\NormalTok{sales\_by\_month}\SpecialCharTok{$}\NormalTok{quarter }\OtherTok{\textless{}{-}} \FunctionTok{ceiling}\NormalTok{(sales\_by\_month}\SpecialCharTok{$}\NormalTok{month }\SpecialCharTok{/} \DecValTok{3}\NormalTok{)}
\CommentTok{\# Create radial plot}
\FunctionTok{ggplot}\NormalTok{(sales\_by\_month, }\FunctionTok{aes}\NormalTok{(}\AttributeTok{x=}\NormalTok{month, }\AttributeTok{y=}\NormalTok{sales, }\AttributeTok{fill=}\FunctionTok{as.factor}\NormalTok{(quarter))) }\SpecialCharTok{+}
    \FunctionTok{geom\_bar}\NormalTok{(}\AttributeTok{stat=}\StringTok{"identity"}\NormalTok{) }\SpecialCharTok{+} \CommentTok{\# use geom\_bar for radial plot}
    \FunctionTok{coord\_polar}\NormalTok{(}\AttributeTok{start=}\DecValTok{0}\NormalTok{) }\SpecialCharTok{+} \CommentTok{\# convert to polar coordinates}
    \FunctionTok{labs}\NormalTok{(}\AttributeTok{title=}\StringTok{"Total Sales by Month"}\NormalTok{, }\AttributeTok{x=}\StringTok{"Month"}\NormalTok{, }\AttributeTok{y=}\StringTok{"Total Sales"}\NormalTok{) }\SpecialCharTok{+}
    \FunctionTok{scale\_fill\_manual}\NormalTok{(}\AttributeTok{values=}\FunctionTok{c}\NormalTok{(}\StringTok{"1"}\OtherTok{=}\StringTok{"lightblue"}\NormalTok{, }\StringTok{"2"}\OtherTok{=}\StringTok{"lightgreen"}\NormalTok{, }\StringTok{"3"}\OtherTok{=}\StringTok{"lightyellow"}\NormalTok{, }\StringTok{"4"}\OtherTok{=}\StringTok{"lightcoral"}\NormalTok{), }\AttributeTok{name=}\StringTok{"Quarter"}\NormalTok{)}
\end{Highlighting}
\end{Shaded}

\includegraphics{djorgensen-520-exercise2_files/figure-latex/radial-1.pdf}

\end{document}
